% !PRINTED 623
\weberpage{623}

%<<CONT
ähnlicher Weise zeigt man, dass keine zwei anderen dieser
Nebengruppen einander gleich sind.

Hiernach ergeben sich für $u$ die innerhalb $A$ conjugirten
Werthe durch Anwendung der Permutationen

\[
(1,\,2)(3,\,4), \quad t(0,1) = (0,1,2,4,3), \quad t(0,2) = (0,2,4,3,1)
\]
\[
t(0,3) = (0,3,1,2,4), \quad t(0,4) = (0,4,3,1,2)\text{:}
\]

\weq{11}{
\begin{aligned}
u_1 &= x_0x_1 + x_1x_2 + x_2x_3 + x_3x_4 + x_4x_0, \\
u_2 &= x_0x_2 + x_2x_1 + x_1x_4 + x_4x_3 + x_3x_0, \\
u_3 &= x_1x_2 + x_2x_4 + x_4x_0 + x_0x_3 + x_3x_1, \\
u_4 &= x_2x_0 + x_0x_4 + x_4x_1 + x_1x_3 + x_3x_2, \\
u_5 &= x_3x_2 + x_2x_4 + x_4x_1 + x_1x_0 + x_0x_3, \\
u_6 &= x_4x_2 + x_2x_0 + x_0x_1 + x_1x_3 + x_3x_4,
\end{aligned}
}

und die entsprechenden $u_1', u_2', u_3', u_4', u_5', u_6'$ findet man, wenn
man auf die $u_1$ die Permutation $t$ anwendet, und dann mit $u_1'$
ebenso verfährt, wie eben mit $u_1$, oder auch einfach, indem man
die in jedem $u$ fehlenden Paare zu dem entsprechenden $u'$ vereinigt
, so dass für alle diese $u$ die Bedingung (7) befriedigt ist.

\weq{12}{
\begin{aligned}
u_1' &= x_0x_2 + x_0x_3 + x_1x_3 + x_1x_4 + x_2x_4, \\
u_2' &= x_0x_1 + x_0x_4 + x_1x_3 + x_2x_3 + x_2x_4, \\
u_3' &= x_0x_1 + x_0x_2 + x_2x_3 + x_3x_4 + x_1x_4, \\
u_4' &= x_0x_1 + x_0x_3 + x_1x_2 + x_2x_4 + x_3x_4, \\
u_5' &= x_0x_2 + x_0x_4 + x_1x_2 + x_1x_3 + x_3x_4, \\
u_6' &= x_0x_3 + x_0x_4 + x_1x_2 + x_1x_4 + x_2x_3.
\end{aligned}
}

Die sechs Grössen (8): $y_1, y_2, y_3, y_4, y_5, y_6$ sind die Wurzeln einer
Gleichung 6\textsuperscript{ten} Grades, deren Coëfficienten rational von $a, b, c, d, e$
und $\sqrt{\varDelta}$ abhängen. Da die $y$ alle ihr Vorzeichen ändern, wenn
$\sqrt{\varDelta}$ das Vorzeichen ändert, so hat diese Gleichung die Form

\weq{12}{
y^6 + a_2y^4 + a_4y^2 + a_6 - \sqrt{\varDelta}(a_1y^5 + a_3y^3 + a_5y) = 0, % [sic] duplicate equation number (12); printed page repeats (12) rather than continuing to (13)
}

worin die $a$ rationale Functionen der Coëfficienten von $f(x)$ sind.
Es müssen aber auch ganze Functionen dieser Coëfficienten sein;
denn $a_2, a_4, \alpha_6, a_1\sqrt{\varDelta}, a_3\sqrt{\varDelta}, a_5\sqrt{\varDelta}$ sind ganze Functionen der % [sic] printed as $\alpha_6$ (lowercase alpha) rather than $a_6$
$x$, die ja hier als unabhängige Variable gelten können, und es
müssen also die drei letzteren durch das Differenzenproduct $\sqrt{\varDelta}$
theilbar sein, und dann müssen sich $a_2, a_4, a_6, a_1, a_3, a_5$ als ganze
symmetrische Functionen der $x$ erweisen. Bestimmt man die
Grade in Bezug auf die $x$, so ergeben sich, da die $y$ vom zweiten
Grade sind, für
%CONT>>
